\documentclass{article}

\usepackage[export]{adjustbox}
\usepackage{listings}
\usepackage{subcaption}
\usepackage{wrapfig}
\usepackage[dvipsnames]{xcolor}
\usepackage{alloy-style}
\usepackage{float}
\usepackage{graphicx}
\usepackage[utf8]{inputenc}
\usepackage[a4paper, top=4cm, bottom=4cm, left=4cm, right=4cm]{geometry}

\title{
    \textbf{\textit{SafeStreets}} \\
    \textbf{DD document}}

\date{Academic year: 2019 - 2020}
\author{
    Dario Miceli Pranio \\
    Pierriccardo Olivieri
}

\begin{document}
\pagenumbering{gobble}

\maketitle

%%%%%%%%%% LOGO POLIMI %%%%%%%%%%
\begin{figure}[h!]
    \centering
    \includegraphics[scale=0.5]{img/logo.png}
\end{figure}

\newpage
\pagenumbering{arabic}
\tableofcontents

\newpage
%%%%%%%%%% CHAPTER 1 %%%%%%%%%%
\section{Introduction}

\subsection{Purpose}
The purpose of this document is to provide a functional description of \textit{SafeStreets} application.

\subsection{Scope}
\textit{SafeStreets} is a service that aims to provide \textit{Users} with the possibility to notify \textit{Authorities} when traffic 
violations occur, and in particular parking violations. The application's goal is achieved by allowing \textit{Users} 
to share photo, position, date, time and type of violation and by enabling \textit{Authorities} to request them.
\\
\\
\textit{SafeStreets} requires the \textit{Users} to create an account to access its services, the functionalities unlocked after 
registration depend on the type of account created.
\\
If a \textit{User} creates an account as \textit{Citizen}, he/she must provide name, surname and a fiscal code in order to prove 
that he/she is a real person. Furthermore, he must provide an email with which he will be uniquely identified 
and a password. Once the account has been activated, \textit{User} can finally start to report parking violations and can also see 
statistics of the streets or the areas with the highest frequency of violations.
\\
\\
On the other hand, an officer will create an account as \textit{Authority} and he will need to provide his name, surname, 
work's Matricola, a password and as for \textit{Citizen}, will be uniquely identified by an email. Once the Matricola 
has been verified and the account has been activated, the officer can retrieve the potential parking violations 
sent by \textit{Citizen} that have not been taken into account yet by other officers, analyze them and, if it is the 
right case, generates traffic tickets. \textit{Authorities}, can see the same statistics of the \textit{Citizen} and can also see
statistics about vehicles' license plate that commit the most violations.


\subsection{Definitions, acronyms, abbreviations}

\subsubsection{Definitions}
\begin{itemize}
    \item \textit{Users}: can be either \textit{Citizen} or \textit{Authority}
    \item \textit{traffic violation}: generic violation that can occur in a street
    \item \textit{parking violation}: a violation caused by a bad parking
    \item \textit{violation}: general violation, identity both traffic or parking violation
    \item \textit{unsafe areas}: areas with an high rate of violations
\end{itemize}

\subsubsection{Acronyms}
Table with all acronyms used in document.
\begin{center}
\begin{tabular}{ | l | l |}
    \hline
    ACRONYM & COMPLETE NAME \\
    \hline
    DD & Design Document \\
    \hline
    GPS & Global Positioning Systems \\
    \hline
    S2B & Software To Be \\
    \hline
    GDPR & General Data Protection Regulation \\
    \hline 
    FC & Fiscal Code \\
    \hline
    UC & Use Case \\
    \hline
\end{tabular}
\end{center}

\subsubsection{Abbreviations}
\begin{itemize}
    \item \textbf{Rn}: n-th Requirement 
\end{itemize}

\subsection{Revision History}

\subsection{Reference documents}
\begin{itemize}
    \item ISO/IEC/IEEE 29148: https://www.iso.org/standard/45171.html
    \item Specification Document: "SafeStreets Mandatory Project Assignement"
\end{itemize}

\subsection{Document Structure}
\begin{itemize}
\end{itemize}
%%%%%%%%%% !CHAPTER 1 %%%%%%%%%%

\clearpage

%%%%%%%%%% CHAPTER 2 %%%%%%%%%%
\section{Architectural Design}
\subsection{Overview}

\subsection{Component view}

\textbf{\textit{SafeStreets} Client Application} \\


\textbf{\textit{SafeStreets} Business Logic} \\
This component describes core logic of \textit{SafeStreets} application. It is a stateless module 
that lies between the Client application and the central Database. Now we will describe each component:
\begin{itemize}
    \item \textbf{Network Manager} \\
    It handles all incoming messages exchanged between Clients and Server. It does two foundamental
    things; it sends all incoming messages from client application to the correct handler in Business
    Logic and on the contrary it collects all the outgoing messages and sends to the corresponding Client
    Application.
    \item \textbf{Authentication Manager} \\
    The Authentication Manager exposes all methods related to the access to the platform. It handles the
    phase of \textit{User} registration and login and in order to do this it has to continuosly interact with 
    the database through the Data Storage Interface. It also checks all the costraints in order to garantee 
    creation of correct account. Furthermore it sends email to all the \text{Users} that have been registered 
    by accessing the Mailing \textit{System}'s Interface. 
    \item \textbf{Citizen Manager} \\
    The \textit{Citizen} manager handles all functionalities related to a \textit{Citizen} Account. In order to
    do this it has to continuosly interact with database through the Data Storage Interface. It also allows
    \textit{Citizen} to update his setting. 
    \item \textbf{Authority Manager} \\
    The \textit{Authority} manager handles all functionalities related to an \textit{Authority} Account. In order to
    do this it has to continuosly interact with database through the Data Storage Interface. It also allows
    \textit{Authority} to update his setting.
    \item \textbf{Data Manager} \\
    This component is the main gateway between the Client Applications and the central Database. It handles the import of 
    new data associated to a \text{User}'s account and its retrieval.
    \item \textbf{Send Report Manager}\\
    The Send Report Manager handles the creation and the management of a Send Report performed by \textit{Citizen} and
    directed to \textit{System}. Inserting is possible thanks to Data Storage Manager. Furthermore, in order to avoid 
    alteration, the Report is encrypted by using Privacy Manager.
    \item \textbf{Retrieve Report Manager}\\
    The Retrieve Report Manager handles the management of a Retrieve Report performed by \textit{Authority}. Retrieving 
    is possible by querying the central Database through the Data Storage Interface.
    \item \textbf{Statistics Manager} \\
    The Statistics Manager handles the management of a Retrieve Statistics performed by \text{User}. Retrieving 
    is possible by querying the central Database through the Data Storage Interface.     
    \item \textbf{Privacy Manager} \\
    The Privacy Manager exposes methods to encrypt sensitive data. This operation is important in order to avoid data leaks
    and data alteration.
    \item \textbf{Data Storage Manager} \\
    The Data Storage Manager provides all methods to interact with the central Database such as data retrieval, 
    storage and update.
    
\end{itemize}

\textbf{External interfaces} \\
Some of the described components in our \textit{System} are also dedicated to communicating with external services through 
specific interfaces. It is essential that the communication works properly in order to fulfill application's functionalities.  

\end{document}